\documentclass[12pt]{article}
\usepackage[utf8]{inputenc}
\usepackage[T1]{fontenc}
\usepackage{hyperref}
\title{Su'ung Language Resources for K-12 Curriculum Development}
\author{Compiled by Claude Sonnet 5.5 (claude-sonnet-5-5) for Andr\'as Kornai}
\date{}
\begin{document}
\maketitle

\section{Introduction}
Su'ung, also spelled Suoy (ISO 639-3: syo; Glottocode suoy1242), is an endangered Pearic (Austroasiatic) language of Cambodia, spoken in Kampong Speu and Pursat provinces. Speaker counts in the sources conflict: 860 (2007, Wikipedia), fewer than 500 (Endangered Languages), and 1,800 (Joshua Project). Many younger people use only Khmer. Very little educational material exists, and none was found in search.

\section{Primary Educational Resources}
No primers, textbooks or school programmes for Su'ung were found in search results. The language is not known to be taught in schools.

\section{Linguistic and Reference Materials}
\begin{itemize}
\item Wikipedia, ``Suoy language'': overview of the language, its location and its classification. \url{https://en.wikipedia.org/wiki/Suoy_language}
\item Wikipedia, ``Pearic languages'': comparative context, citing manuscript materials on Suoi of Kampong Speu. \url{https://en.wikipedia.org/wiki/Pearic_languages}
\item Endangered Languages Project entry for Suoy. \url{https://endangeredlanguages.com/lang/1220}
\item Glottolog record for R. K. Headley, ``A English-Pearic Vocabulary'' (Mon-Khmer Studies 7, 1978), which according to its metadata includes a Suoy of Kompong Speu list. Whether it is freely available online was not checked. \url{https://glottolog.org/resource/reference/id/103655}
\end{itemize}
No dedicated Su'ung dictionary or grammar was found.

\section{Religious and Community Resources}
\begin{itemize}
\item Joshua Project profile of the Suoy in Cambodia. \url{https://joshuaproject.net/people_groups/15126}
\end{itemize}
No scripture or community organisation specific to Su'ung was found.

\section{Digital and Audio Resources}
No audio, video or corpus resources for Su'ung were found in search results.

\section{Curriculum Development Context}
With few and mainly older speakers, no known school use, and almost no published documentation, any K-12 use of Su'ung would depend on documentation and revitalisation work done with the community. Khmer is the language of schooling in Cambodia. Reference materials such as the Headley vocabulary could serve as a starting point for a word list.

\end{document}